% 1. Pass the hyphens option BEFORE loading the class
\PassOptionsToPackage{hyphens}{url}

\documentclass[sigconf]{acmart}
\setcitestyle{numbers,sort&compress}

% Student project formatting
\setcopyright{none}
\settopmatter{printacmref=false}

% 2. Add the hyphenat package with the 'htt' option 
% (This allows texttt blocks to break across lines)
\usepackage[htt]{hyphenat}
\usepackage{xurl} % Keeps your bibliography links wrapped

% Student project formatting
\setcopyright{none}
\settopmatter{printacmref=false}
\renewcommand\footnotetextcopyrightpermission[1]{}
\pagestyle{plain}
\usepackage{booktabs}
\usepackage{graphicx}
\usepackage{url}
\def\UrlBreaks{\do\/\do-\do\_}

\begin{document}

\title{Interactive Exploratory Search of Scholarly Literature Using Multi-View Visual Analytics}

\author{Mosfaiul Telot}
\affiliation{%
  \institution{University of Regina}
  \city{Regina}
  \state{Saskatchewan}
  \country{Canada}
}
\email{mosfaiul.telot@uregina.ca}

\author{Trace Price}
\affiliation{%
  \institution{University of Regina}
  \city{Regina}
  \state{Saskatchewan}
  \country{Canada}
}
\email{price23t@uregina.ca}

\author{Nilab Naderi}
\affiliation{%
  \institution{University of Regina}
  \city{Regina}
  \state{Saskatchewan}
  \country{Canada}
}
\email{nilab.naderi@uregina.ca}

\begin{abstract}
Exploratory search in scientific literature is difficult when users rely on ranked lists that provide limited support for understanding topic structure, temporal change, and paper-to-paper relationships. We present a web-based visual analytics system that supports exploratory literature search through coordinated views of papers, terminology, and authors. The system uses OpenAlex as the source, a Node.js/Express backend for retrieval and aggregation, MongoDB for storage, and a React frontend for interactive analysis. Our workflow includes retrieval filtering for article and proceedings-article records, terminology aggregation from multiple metadata fields, and linked interactions across dashboard views. A scenario-based evaluation indicates that the system helps users identify influential papers, detect emerging terms, and move between overview and detail more effectively than list-only workflows. We discuss limitations related to metadata quality, source coverage lag for recent publications, and scalability constraints for larger retrieval sets.
\end{abstract}

\begin{CCSXML}
<ccs2012>
<concept>
<concept_id>10002951.10003260.10003261</concept_id>
<concept_desc>Information systems~Search interfaces</concept_desc>
<concept_significance>500</concept_significance>
</concept>
<concept>
<concept_id>10002951.10003260.10003282</concept_id>
<concept_desc>Information systems~Search interaction</concept_desc>
<concept_significance>300</concept_significance>
</concept>
<concept>
<concept_id>10003120.10003121.10003124</concept_id>
<concept_desc>Human-centered computing~Visualization systems and tools</concept_desc>
<concept_significance>300</concept_significance>
</concept>
</ccs2012>
\end{CCSXML}

\ccsdesc[500]{Information systems~Search interfaces}
\ccsdesc[300]{Information systems~Search interaction}
\ccsdesc[300]{Human-centered computing~Visualization systems and tools}

\keywords{exploratory search, literature visualization, visual analytics, bibliometric analysis, OpenAlex}

\maketitle

\section{Introduction}
The volume of scholarly publications continues to grow rapidly. Conducting a literature review therefore remains slow and cognitively demanding: millions of works exist, ranked lists and abstract skimming offer little structure, and it is easy to feel overwhelmed when choosing what to read next. Researchers often start from a single lead, such as a topic phrase or keyword, yet need a \emph{curated} slice of the literature and a sense of how ideas connect before committing hours to full-text reading.

List-based search is strong for known-item retrieval but weak for this exploratory phase: it hides clusters of related work, branching lines of inquiry, and opportunities to combine concepts (for example, linking exploratory search with metacognition or large language models). This challenge motivates a visual analytics approach that surfaces relationships, ranked terminology, and user-driven filtering so that analysts can \emph{preattentively} narrow promising reading paths.

This project investigates the question: \textit{How can a coordinated multi-view visualization system support literature-review exploration when the user begins with only a topic lead and a broad domain?} We implement an end-to-end system that integrates OpenAlex retrieval, backend processing and aggregation, and interactive frontend visualizations for papers, terminology, and authors.

Our contributions are:
\begin{itemize}
  \item A practical architecture for exploratory literature analysis using OpenAlex, Node.js/Express, MongoDB, and React.
  \item A coordinated dashboard combining ranked tables, relationship graphs, and trend visualizations.
  \item A terminology aggregation approach that unifies primary topics, tags, and keywords for trend and relationship analysis.
  \item Favorite-term filtering with match-to-confidence scoring and linked table--graph highlighting to support synthesis-oriented exploration.
  \item A scenario-based demonstration of exploratory workflow from overview to detailed inspection.
\end{itemize}

\section{Problem Domain and User Tasks}

\subsection{Problem statement}
Anyone who has started a literature review knows the feeling: there are too many papers, and the clock is ticking. You open tab after tab, skim abstracts, and still wonder whether you are missing an entire thread of work or a clean way to combine two ideas you care about. The central problem we address is that cost---not storage, but \emph{orientation}. With vast numbers of publications, it is easy to spend many hours and still lack a clear sense of \emph{where} active sub-communities sit, \emph{how} a theme has forked, and \emph{which} combinations of concepts might become a fresh research angle.

Most search interfaces still push long text lists. Lists hide structure. They do not show clusters, branches, or repeated vocabulary that might nudge you toward a first reading list or a mashup across subfields.

Our system targets that early, exploratory phase. The user often has only a small lead. For example a short query such as ``exploratory search.'' They want an honest snapshot of what is happening in the literature, but not the whole internet. We therefore retrieve a \emph{domain-scoped}, connected sample (not the entire OpenAlex index) and present it through linked views: papers, terms, citations, and authors. Optional filtering helps users pick promising paths before they invest in full-text reading.

\subsection{Target users}
Primary users are graduate students, researchers, and advisors who must survey a topic, identify influential or representative papers, and discover adjacent concepts. A typical motivation is to shorten the path from ``I have a keyword'' to ``I have a shortlist and a reading order.''

\subsection{User tasks and how the interface supports them}
We decompose the literature-review exploration workflow into concrete tasks supported by the dashboard:

\begin{enumerate}
  \item \textbf{Seed the exploration with a topic and domain.} The user enters a query on the home page and selects a domain (dropdown mapped to OpenAlex field/subfield filters). The backend retrieves a curated, connected set of articles and proceedings articles. This addresses the need to bound the search space when the user's lead is underspecified.

  \item \textbf{Understand structure at a glance via the paper network.} The force-directed tree links papers by citation and terminology overlap. Tight \emph{clusters} often read as ``people working in the same neighborhood.'' A hub with many children can read as ``this idea spawned a lot of follow-on work.'' You still have to read eventually, but color and layout give you a cheap first pass: where to click first without opening every abstract.

  \item \textbf{Mine vocabulary for new directions.} The ranked terminology table aggregates keywords, tags, and primary topics. Scanning these terms can trigger combinatorial hypotheses (e.g., combining exploratory search with metacognition, or with LLM-mediated interaction) that a plain title list might not surface.

  \item \textbf{Curate and combine ideas with favorite terms and filtering.} The user pins terms in a favorites panel at the bottom of the dashboard, then hits \emph{Filter}. The paper table keeps only works that share at least one of those terms. Papers that match more of the chosen terms score higher on a simple \emph{confidence} bar. The graph uses the same score for node color. The point is practical: you can ``pre-attentively process'' which papers align with your conceptual intersection before you commit to a reading order.

  \item \textbf{Follow reading paths without exhaustive full-text pass.} Citation and term edges, together with filtered highlights, suggest natural next papers (e.g., heavily cited neighbors or papers sharing multiple favorite concepts). The intent is to substitute hours of undirected abstract reading with a smaller set of structurally justified entry points for the actual literature review.

  \item \textbf{Broaden or deepen with authors and trends.} Author aggregates and topic timelines complement the network: users can shift from paper--paper structure to who publishes where and how term prevalence changes over time, supporting decisions about depth versus breadth.
\end{enumerate}

\subsection{Design goal}
We want to lower the mental load of the first pass through a new topic. Retrieval alone is not enough if the user still has to glue pieces together in their head.

So we tie views together on purpose. The paper table, the force graph, the ranked terms, and the favorites filter should behave like one workspace. Structure and vocabulary should nudge the user. Their own pinned terms should steer what highlights and what fades.

The end goal is modest but useful: help someone leave the dashboard with a short, defensible list of papers to actually read---not with perfect coverage, but with a clearer story about where to start.

\section{Literature Review}

When relying on traditional search engines to locate academic information, researchers are frequently confronted with flat, ranked lists of documents. While effective for known-item retrieval, this paradigm imposes a high cognitive load during exploratory search, making it difficult to understand the structural relationships and temporal evolution of large research collections \cite{hearst2009search, marchionini2006exploratory}. To solve this problem, previous studies have heavily focused on applying visualization techniques to assist in the analysis of scientific literature, aiming to visually represent trends, relationships, and significant contributions. 

A major contribution to this field is CiteSpace, developed by Chen (2006) \cite{chen2006citespace}. Chen argues that citation information can be utilized to determine significant papers and promising research areas over time. Through the analysis of citation frequencies and topic evolution, users gain a clearer understanding of how a research domain matures. This approach confirms that research activity is highly dynamic and best investigated using temporal visualizations. Alongside determining historical trends, comprehending the spatial relationship between topics is equally critical. VOSviewer, developed by Van Eck and Waltman \cite{vaneck2010vosviewer}, is a widely adopted software for building and visualizing bibliometric networks. Their work demonstrates how relationships among topics, authors, and publications can be mapped using network graphs, where strongly related items are pulled closer together. Such visualizations allow users to identify distinct research clusters, which is particularly beneficial when navigating interdisciplinary fields.

Further studies on mapping scientific knowledge confirm the utility of visualizing entire research domains. Börner (2010) \cite{borner2010atlas} describes the application of science mapping methods to uncover the organization and development of knowledge fields, allowing users to better comprehend how scientific knowledge evolves. Furthermore, topic-based visualization systems like TIARA \cite{liu2012tiara} demonstrate how large document collections can be tracked over time. These systems enable users to explore topics dynamically, understanding how particular themes rise, expand, or disappear, proving that research topics are not fixed but constantly fluctuating. 

The theoretical foundation for these systems lies in the broader field of visual analytics. Keim et al. (2010) \cite{keim2010visualanalytics} define visual analytics as the unification of automated data analysis with interactive visualizations, enabling human sensemaking of complex, massive datasets. This highlights the critical importance of human-in-the-loop interaction; users must be able to engage with data, make decisions, and narrow down their knowledge via visual interfaces.

While traditional tools like VOSviewer and CiteSpace excel at mapping science, they are primarily designed for static, retrospective scientometric analysis of manually exported datasets. Modern exploratory search requires real-time responsiveness. Building upon Keim’s visual analytics framework and Börner’s science mapping, our proposed project aims to bridge this gap. By creating an interactive visualization platform that synthesizes real-time semantic fetching with multi-view visualizations (timeline, citation, and relational force-directed graphs), this system actively supports exploratory search and dynamic literature discovery.

Beyond bibliometric mapping, the theoretical foundation of this work is rooted in exploratory search and sensemaking. Exploratory search represents an iterative process of learning and investigation \cite{marchionini2006exploratory, white2009exploratory}, where researchers must manage high cognitive load by externalizing their evolving understanding into manipulable representations \cite{pirolli2005sensemaking, kirsh2010thinking}. Traditional search interfaces fail this process by relying on flat, ranked lists \cite{hearst2009search}. To address this, our system employs coordinated and multiple views (CMV) \cite{roberts2007state} and follows Shneiderman’s Visual Information-Seeking Mantra: "Overview first, zoom and filter, then details-on-demand" \cite{shneiderman1996eyes}. By leveraging preattentive visual variables such as color and spatial position \cite{ware2019information}, and adhering to effective visualization design principles \cite{munzner2014visualization}, the system allows researchers to recognize patterns and "mashup" ideas through faceted keyword interaction \cite{kules2009exploratory} and semantic similarity scoring \cite{ammar2018construction, wang2020microsoft}.

\section{System Overview and Architecture}

The system architecture is designed to support real-time exploratory search and is divided into three primary tiers: Data Orchestration (Retrieval and Storage), the Semantic Analytical Engine (Backend Processing), and the Coordinated Visualization Layer (Frontend UI).

\subsection{Data Orchestration: The Seed Paper and Snowball Fetch}
The data pipeline is initiated via a custom fetching service (\texttt{fetchDa\allowbreak ta.js}) that interfaces with the OpenAlex API. The process begins when a user's initial search query retrieves a highly relevant "seed paper" (or primary set of seed works). Rather than presenting just this initial keyword match, the system employs a "snowball fetch" methodology. It uses the seed paper as an anchor to recursively retrieve its entire network: its direct citations, referenced works, associated authors, and underlying topic tags. To ensure data quality and visual consistency, the backend normalizes metadata—such as converting specialized characters and formatting mathematical notation via MathJax—before persisting the structured documents into a MongoDB database. This database acts as the single source of truth, caching historical citation counts, author arrays, and terminology tags.

\subsection{Semantic Analytical Engine: Beyond Keyword Matching}
The backend, built on Node.js and Express, acts as the intelligence layer of the application. It goes beyond standard CRUD operations by actively computing relationships between the cached papers to determine the most valuable suggestions.
\begin{itemize}
    \item \textbf{Terminology and Trending Services:} The \texttt{terminologyService.js} module aggregates keywords across the dataset, applying a custom stop-word blacklist to filter out generic academic terms. It calculates temporal "trend scores" based on citation frequency over a 10-year rolling window.
    \item \textbf{Network Construction and Similarity Scoring:} The \texttt{paperNetworkService.js} dynamically builds a directed graph of the literature. To determine which papers are most relevant to the user's exploration, the system computes a custom semantic similarity score between the seed paper and all other fetched papers. This algorithm mathematically weights shared variables: common authors, overlapping concepts, shared keywords, and direct citations.
    \item \textbf{Superior Suggestion Quality:} This multi-dimensional semantic scoring dictates the edges (links) and clustering behavior passed to the frontend. It provides vastly superior suggestions compared to traditional search engines because it overcomes the common "vocabulary mismatch" problem. Even if a related paper uses completely different terminology in its abstract, it will achieve a high similarity score if it shares the same foundational citations and concepts as the seed paper. This allows the system to accurately surface interdisciplinary connections and hidden research pathways that a manual keyword search would miss.
\end{itemize}

\subsection{Coordinated Visualization Layer}
The frontend interface, developed in React, orchestrates a multi-view visual analytics dashboard. It consumes the REST endpoints provided by the backend to render a highly interactive workspace. The interface relies heavily on state management (\texttt{App.js}) to conditionally render different dashboards based on the user's level of exploration (e.g., Home Search, Main Dashboard, or Single Paper/Author drill-downs). 

The visual layer integrates two primary graphics libraries to handle complex data rendering:
\begin{itemize}
    \item \textbf{amCharts 5:} Powers the \texttt{PaperForceGraph.jsx}, generating a highly interactive, force-directed hierarchy. This module handles real-time node colorization based on user-selected confidence scores and keyword filtering.
    \item \textbf{D3.js:} Drives the coordinated supplementary views, including the \texttt{AuthorBarChartN.jsx} for ranking researcher impact, \texttt{CitedLineChart.jsx} for tracking temporal citation histories, and \texttt{WordCloud.jsx} for spatial terminology exploration. 
\end{itemize}
Through bidirectional state synchronization (\texttt{BackFetch.js}), user interactions in one view (such as adding a term to the "Favorite Keywords" panel) dynamically update the backend queries, instantly filtering and re-rendering the visual components across the entire dashboard.

\section{Backend and Implementation}
\subsection{Server stack and configuration}
The backend is a Node.js application using Express. It enables CORS for the React development server, parses JSON request bodies, and connects to MongoDB through Mongoose using a connection string from environment variables (\texttt{MONGO\_URI}). Optional \texttt{USER\_EMAIL} is passed to OpenAlex API requests as the recommended \texttt{mailto} parameter for polite pooling. The server listens on a configurable port (default 5000).

\subsection{Data ingestion from OpenAlex}
Ingestion is implemented in the fetch module and triggered by \texttt{POST /api/trigger-fetch} with a search string and optional category (mapped to OpenAlex topic field filters, e.g.\ computer science subfields). The pipeline: (1)~query OpenAlex for a small set of seed works matching the search; (2)~fetch works that cite those seeds; (3)~fetch referenced works (chunked to respect URL limits); (4)~merge and deduplicate by OpenAlex work ID. Requests use HTTPS with retry and exponential backoff on HTTP 429 (rate limit). Each run clears the \texttt{Paper} collection and inserts a fresh corpus so the dashboard reflects the latest query.

Works are filtered at the API and again before save: only \texttt{article} and \texttt{proceedings-article} types are kept, with English language and optional subfield/topic filters. Basic quality checks require a non-trivial title and at least one author. Stored fields include OpenAlex ID and URL, DOI, title, year, citation count, venue, reconstructed abstract, author list, keywords, primary topic, tags (from topics and higher-level concepts), referenced work IDs, per-year citation counts, and work type.

\subsection{MongoDB schema and auxiliary collections}
Scholarly metadata lives in a Mongoose \texttt{Paper} model (document-shaped records). User-specific preferences use separate MongoDB collections: \texttt{favorite-papers} and \texttt{favorite-terms}, keyed by search term, so favorites persist across sessions without altering core paper documents.

\subsection{REST API surface}
Table~\ref{tab:api} summarizes major routes the dashboard and auxiliary pages use. Aggregation-heavy logic runs in route handlers or imported services (\texttt{paperNetworkService}, terminology helpers). The table uses the two-column \texttt{table*} environment so it spans the full page width under the ACM \texttt{sigconf} template (avoids overflow into the neighboring column).

\begin{table*}[t]
  \centering
  \small
  \caption{Representative backend HTTP endpoints.}
  \label{tab:api}
  \begin{tabular}{@{}p{0.30\textwidth}p{0.64\textwidth}@{}}
    \toprule
    Method \& path & Role \\
    \midrule
    \texttt{POST /api/trigger-fetch} & Start OpenAlex ingestion for a query \\
    \texttt{GET /api/fetch-status} & Poll ingestion progress \\
    \texttt{GET /api/papers/dashboard-stats} & Aggregate counts, top tag, sample top papers \\
    \texttt{GET /api/top-cited} & Ranked papers (optional \texttt{workType} filter) \\
    \texttt{GET /api/terminology} & Ranked terms with trend series \\
    \texttt{GET /api/topic-timeline} & Concept counts per year for timelines \\
    \texttt{GET /api/paper-network} & Nodes and citation/terminology edges for graphs \\
    \texttt{GET /api/authors}, \texttt{/api/tags} & Author and tag aggregates \\
    \texttt{GET/PUT /api/favorite-paper/:term} & Load/save favorited paper IDs \\
    \texttt{GET/PUT /api/favorite-term/:term} & Load/save favorited keywords \\
    \texttt{GET /api/closest-papers/:id} & Similarity-ranked related papers \\
    \bottomrule
  \end{tabular}
\end{table*}

\subsection{Paper network and terminology computation}
\texttt{GET /api/paper-network} delegates to a service that loads papers from MongoDB and constructs nodes (internal IDs, citation counts, term fields) and edges: pairwise terminology overlap (shared tags/keywords/topics) and citation edges when referenced OpenAlex IDs resolve to ingested papers. This supports both citation-mode and terminology-mode graph layouts on the client.

Ranked terminology (\texttt{/api/terminology}) aggregates tags, keywords, and primary topic across the corpus, applies a small stoplist for generic labels, and computes counts and recent-year trend vectors for sorting (e.g.\ by frequency or recency).

\subsection{Frontend--backend contract}
The React app targets \texttt{http://localhost:5000} for API calls during development (with a proxy option for same-origin requests). The main dashboard polls dashboard stats until papers exist, then batches requests for timeline, terminology, paper network, top-cited list, tags, and authors. This keeps the UI responsive while ingestion may still be running.


\section{Visualization Design}
The main page uses a stratified dashboard layout: the first row emphasizes paper relationships, the second row emphasizes term trends, and the third row emphasizes author information. This layout follows abstraction and simplicity by showing only the most decision-relevant information in each region. We intentionally keep background ink, decorative borders, and non-data strokes minimal so visual emphasis remains on the data marks (high data-to-ink ratio).

\subsection{Design principles applied from class}
\textbf{Gestalt closure/common region.} The top-papers area groups the table and the force tree into one enclosed region because both represent the same paper set from different perspectives. This supports chunking: users read them as one analytical unit instead of two unrelated widgets.

\textbf{Gestalt connectedness.} In the force tree, links explicitly encode relationships between papers (citation or terminology). Connected nodes are perceived as related work streams, making cluster and bridge structure visible at a glance.

\textbf{Detail-on-demand.} The dashboard follows "overview first, zoom/filter, details on demand": users can hover nodes and edges for local details, and click papers/authors for deeper drilldown views, while the default state remains visually simple.

\textbf{Visual channels and McKinlay's ranking.} We use \emph{size} to encode citation magnitude on graph nodes. For ranked keywords (nominal categories), \emph{color hue} is used as the primary channel. For confidence/strength (ordinal and continuous), we use a red-to-green continuous color ramp; this communicates ordered intensity rather than discrete classes. In addition, color saturation is used as a secondary cue for citation/importance context where appropriate.

\textbf{Preattentive consistency and change blindness mitigation.} The same color mapping is reused across linked views (table bars and graph nodes), so users can preattentively trace high-confidence papers across components. The fixed favorite-keyword panel at the bottom preserves context while other views update, reducing change blindness during filtering.

\textbf{Opponent-process color channel.} The confidence palette leverages a red-green opponent channel to convey low-to-high match strength. In our design, red indicates weak alignment and green indicates strong alignment with selected favorite terms.

\subsection{Force Paper Connected Graph}
The main interactive element in the dashboard is the paper force-directed graph. In the force directed graph each node is represented by a node.  This is a graph which has two connection modes: one for citations and one for terminology. In the different modes the connections between the nodes change. Citations of each paper determine the size of each node. 

\begin{figure}[H]
    \centering
    \includegraphics[width=1.00\linewidth]{MainGraph.png}
    \caption{Main Graph}
    \label{fig:enter-label}
\end{figure}

The graph is positioned on the right and the corresponding papers are positioned on the right side. The nodes can be moved around by the user allowing them to clearly see all of the edges. The weight of each node is based on the number of edges that are connected to it.  Initially the nodes are grey, colour is used when the nodes are filtered out, so the nodes are not initially given unique colours. The edges are given colours to indicate clusters.  Central nodes have a circle around them to indicate that the node has several neighbors. The edge and radius colouring show the user which nodes are indirectly related and which nodes are not. 

\begin{figure}[H]
    \centering
    \includegraphics[width=1.00\linewidth]{filter.png}
    \caption{Main Graph}
    \label{fig:enter-label}
\end{figure}

In order to personalize the paper force graph we allow the user to select topics and filter the graph based on those topics. When the user selects the heart beside the paper name all of the topics associated with that paper will be added to the favorite topics list. The user can also add individual topics from the topic ranked table. The user can freely modify the favorite topics by adding and removing them from the favorite topics page. When the user filters the topics only the papers that relate the most with those topics will be shown. The papers will be ranked based on how close they are to the topic and given a colour between red and green. The user can see in the graph, which papers most closely resemble their filters. The scale below the paper names show the colour scale. A red-green colour scale is used because it is a colour scale the user are most familier with due to its uses in stock market tickers or traffic lights. 

To allow the user to quickly understand the trends of various keywords we used a scrollable mini-chart. 

\begin{figure}[H]
    \centering
    \includegraphics[width=1.00\linewidth]{Topics.png}
    \caption{Here is a scrollable mini-chart alongside a line chart}
    \label{fig:enter-label}
\end{figure}

The mini-chart allows users to efficiently search through all of the key-words related to the topic that they searched. The minichart is used because the user needs to only see the trends and not the exact values, which are unneeded. The minichart can allow the user to find other fields that they are interested in to quickly move to. The user also has an option to drilldown by clicking on one of the key words and the chart for that key-word appears on the main line graph. In the main line graph the user can hover over the lines with their mouse to get the exact values. Colours are given to each keyword due to the amount of keywords they are not each given a unique colour. The colour in the minicharts alwayse matches the colour in the line graph.  

At the bottom of the dashboard we have a bar chart to show the most prolific authors in the field. The user can quickly read through the list of authors and compare their number of papers in the field that the user search or the number of citations in those papers. 

\begin{figure}[H]
    \centering
    \includegraphics[width=1.0\linewidth]{AuthorBarChart.png}
    \caption{Bar Chart}
    \label{fig:enter-label}
\end{figure}

In the author bar chart we use distinct colours with adequate spacing to ensure that the user can tell apart the bars for different authors. The authors are sorted by the number of papers that they released because the number of papers show how often an author published in that field. Since there are two bars for each author we use a key to tell them apart. Since the bars can be far apart from the axis we allow the user to hover over the bar to get the exact value.  With frequent use a user is able to remember two colours without looking at the key. We made the names of the authors blue rather than gray, to ensure that the user understands that they have the option to click on the labels.

To allow the user to get more information about a given paper or author we provide two miniature dashboards one for the paper and one for the author. 

\begin{figure}[H]
    \centering
    \includegraphics[width=1.0\linewidth]{pDash.png}
    \caption{Paper Dashboard}
    \label{fig:enter-label}
\end{figure}

The paper dashboard is focused on giving information about the paper to the use. Since there is limited data in a single paper much of the information presented on the paper dashboard is direct data. Since there is a relativly small amount of authors, it is easiest to give the authors information directly to the user and no using a visualization. We also give the user the abstract for the paper, giving them more detailed information about the contents of the paper. The also show the users the main topics of the paper and allow them to add those topics to their favorite topics list. 
We want to allow users to select a paper and be able to find similar papers in the field. We calculated a score using concepts, coauthors and common citations to determine what the closest papers are to the selected papers. We presented them as a list rather than a visualization, because the order of closest papers matter to the user more than the score of the papers.
The only visualization in the single paper dashboard is a line chart depicting the citations per year of the paper. The line chart is used because the user in most interested in citation trends and the angle of the line chart can be pre-attentivly processed. This means the user can quickly understand the trend of the line chart with out comparing the positions of the line segments.

The author dashboard page is similar to the paper dashboard page in the content. It also has a line graph showing citations per year. The dashboard lists the most commen co-authors and the authors papers. The main use of the author dashboard is the allow the user to find more papers from a selected author.

The interface also supports explicit drilldown: clicking a paper in the ranked table opens the paper-level dashboard, and clicking an author in the author bar chart opens the author-level dashboard. This keeps navigation consistent with exploratory workflows where users repeatedly move from overview to evidence and back.




\section{Uses Cases}
In order to illustrate the usefulness of the system, a sample scenario is taken into account. The usefullness of the system is that users can quickly search up different concepts keys sequenctially and get useful information about that topic quickly. This allows a new student unsure of what they want to study to quickly look through a bunch of different fields and to find the top authors in those areas. In any given search if the concept graph is typed into the system. The system retrieves pertinent research papers and tracks through them to find there connections based on common citations and concepts. This chart shows the relations between papers in the field of graphs based on their core ideas. The user can drill down on a topic and get a graph with connections based on that topic.  The line chart illustrates the development of research on graphs over the years. The line chart will also go through all of the related concepts to graphs such as Advanced Graph Neural Networks or Data Mining to show the user the trends on all of these topics. The user can scroll down to the bar chart to find the authors in the field and look at the author page. In the author page the can see all of the collected papers that that author wrote and their co-authors. The user can drill down on a given paper to see the authors of the paper and the closest papers that the search collected.  The network graph shows the correlation among various papers and concepts. These visualizations allow the user to grasp the structure of the research field in a short time without having to read individual papers.

\section{Benefits and Limitations}  
The system suggested has a number of significant benefits to users of scientific literature. First, it enhances clarity and understanding greatly by converting extensive textual data into visual informations. Users do not need to read long lists of papers to understand trends, patterns and relationships; through charts and graphs they can quickly interpret these relationships, patterns and trends. This streamlines the research process and makes it more intuitive. Second, the system assists in lessening the cognitive burden on users. In the case of traditional search engines, users will have to manually compare and analyze various documents, which can be time-consuming and mentally taxing. Conversely, this system displays summarized information in the form of visualizations, which enables users to comprehend advanced information with reduced efforts. In this system 
Interactive exploration is also another great advantage. Users do not need to view the results as being static but they can interactively explore various aspects of the data. As an illustration, they will be able to see the change of topics over time with the help of line charts, find the most prolific authors in the field with the help of bar charts, and see connections between papers based on their citations and their papers with networks graphs. This interactive style encourages searching and browsing, in which users do not necessarily want to find a predetermined purpose but desire to make discoveries. The system also facilitates trend analysis and decision making. Users can see how research activity is changing over time, and how it is possible to identify emerging topics, areas of research popularity, and those that are on the downward trend. It is especially helpful to students who choose the topics of their research or researchers who study the gaps in the literature. In spite of these positive features, there are a number of limitations to the system. A major constraint is that it is dependent on data quality and availability. The accuracy and usefulness of the visualizations depend on the quality of the data retrieved from external sources such as OpenAlex. When the data is not complete or divergent, then the outcomes can be misleading. Unfortunately, there is coverage gaps in the OpenAlex's API so there is missing data in recent years. The other issue with the data is the labeling of concepts. In some cases the API label a paper as having a certain concept which does not accurately describe the topics that the paper covered. Performance and scalability is another constraint. The system can also slow in processing time when dealing with large datasets particularly in fetching data from the backend and visualization rendering. This may impact user experience, especially when you need a great number of papers.  The system is also complicated. It is based on several technologies such as React, Node.js, MongoDB and D3.js. This makes development and maintenance more complex, and needs technical knowledge to maintain and enhance the system.
While the system effectively supports exploratory research, evaluating the force-directed graph imposes a preliminary learning curve. The interface, though designed for usability, requires a brief onboarding phase for users to become proficient with its multi-dimensional filtering capabilities. Currently, the platform focuses on macro-level topic trends and relationship mapping. Future iterations will aim to integrate deeper semantic analysis, personalized recommendation algorithms, and real-time data synchronization to further reduce user cognitive load.

\section{Conclusion}
In conclusion, this project presents an interactive, multi-view visual analytics platform designed to advance exploratory search in scholarly literature. By synthesizing modern backend infrastructure with robust data visualization libraries, the system successfully transforms static bibliographic data into manipulable, relational networks. This approach addresses the fundamental limitations of traditional, list-based search interfaces by enabling researchers to rapidly evaluate conceptual intersections, map domain trends, and externalize their sensemaking process. Future development will focus on optimizing rendering performance for larger datasets and expanding the combinatorial filtering capabilities of the visual workspace.

\bibliographystyle{ACM-Reference-Format}
\bibliography{references}

\end{document}